
\documentclass{article}
\usepackage{colortbl}
\usepackage{makecell}
\usepackage{multirow}
\usepackage{supertabular}

\begin{document}

\newcounter{utterance}

\centering \large Interaction Transcript for game `hot\_air\_balloon', experiment `air\_balloon\_survival\_lv\_complexity\_easy', episode 1 with Qwen3.6{-}35B{-}A3B{-}FP8{-}without{-}reasoning.
\vspace{24pt}

{ \footnotesize  \setcounter{utterance}{1}
\setlength{\tabcolsep}{0pt}
\begin{supertabular}{c@{$\;$}|p{.15\linewidth}@{}p{.15\linewidth}p{.15\linewidth}p{.15\linewidth}p{.15\linewidth}p{.15\linewidth}}
    \# & $\;$A & \multicolumn{4}{c}{Game Master} & $\;\:$B\\
    \hline

    \theutterance \stepcounter{utterance}  
    & & \multicolumn{4}{p{0.6\linewidth}}{
        \cellcolor[rgb]{0.9,0.9,0.9}{
            \makecell[{{p{\linewidth}}}]{
                \texttt{\tiny{[P1$\langle$GM]}}
                \texttt{Jūs piedalāties kopīgā sarunu spēlē.} \\
\\ 
\texttt{Kopā ar citu dalībnieku jums jāvienojas par vienu priekšmetu kopumu, kas tiks saglabāts. Katram no jums ir savs viedoklis par to, cik liela nozīme jums ir katram priekšmetam (nozīmīgums). Jūs nezināt, kā otrs dalībnieks vērtē priekšmetus. Turklāt jums ir norādīts, cik pūļu prasa katrs priekšmets.} \\
\texttt{Jūs varat vienoties par kopumu tikai tad, ja izvēlēto priekšmetu kopējais piepūles apjoms nepārsniedz kopīgo robežu:} \\
\\ 
\texttt{LIMIT = 5492} \\
\\ 
\texttt{Šeit ir norādītas atsevišķu priekšmetu piepūles vērtības:} \\
\\ 
\texttt{Item effort = \{"C08": 106, "C55": 446, "A43": 933, "A67": 988, "B42": 388, "A62": 826, "B19": 994, "C43": 556, "C73": 932, "C10": 838, "B36": 302, "A18": 564, "C61": 260, "C67": 729, "C42": 489, "C70": 323, "B40": 103, "C39": 213, "A21": 668, "C16": 326\}} \\
\\ 
\texttt{Šeit ir jūsu personīgais viedoklis par katra priekšmeta nozīmīgumu:} \\
\\ 
\texttt{Item importance values = \{"C08": 138, "C55": 583, "A43": 868, "A67": 822, "B42": 783, "A62": 65, "B19": 262, "C43": 121, "C73": 508, "C10": 780, "B36": 461, "A18": 484, "C61": 668, "C67": 389, "C42": 808, "C70": 215, "B40": 97, "C39": 500, "A21": 30, "C16": 915\}} \\
\\ 
\texttt{Mērķis:} \\
\\ 
\texttt{Jūsu mērķis ir vienoties par tādu kopīgu priekšmetu kopumu, no kura Jūs gūsiet pēc iespējas lielāku labumu (t.i., kas maksimāli palielina kopējo nozīmīgumu JUMS), vienlaikus nepārsniedzot LIMIT. Katrā ziņojumā jums nav obligāti jāiesniedz PRIEKŠLIKUMS {-} jūs varat vienkārši arī sarunāties. Visas taktikas ir atļautas!} \\
\\ 
\texttt{Saskarsmes protokols:} \\
\\ 
\texttt{Ziņojumā drīkst izmantot tikai šādus strukturētus formātus:} \\
\\ 
\texttt{PRIEKŠLIKUMS: \{'A', 'B', 'C', …\}} \\
\texttt{Ierosiniet saglabāt tieši šos priekšmetus.} \\
\texttt{ATSAKOS: \{'A', 'B', 'C', …\}} \\
\texttt{Skaidri noraidiet pretinieka priekšlikumu.} \\
\texttt{ARGUMENTS: \{'...'\}} \\
\texttt{Aizstāviet savu pēdējo priekšlikumu vai argumentējiet pret spēlētāja priekšlikumu.} \\
\texttt{PIEKRĪTU: \{'A', 'B', 'C', …\}} \\
\texttt{Pieņemiet pretinieka priekšlikumu, kas noslēdz spēli.} \\
\\ 
\\ 
\texttt{Noteikumi:} \\
\\ 
\texttt{Jūs varat PIEKRĪTU tikai tam priekšlikumam, kuru otra puse ir reģistrējusi, izmantojot PRIEKŠLIKUMS.} \\
\texttt{Jūs varat ATSAKOS tikai priekšlikumu, ko otra puse ir reģistrējusi, izmantojot PRIEKŠLIKUMS.} \\
\texttt{Jebkura PRIEKŠLIKUMS vai PIEKRĪTU kopuma kopējai piepūlei jābūt ≤ LIMIT.} \\
\texttt{NEATKLĀJIET savus slēptos nozīmīguma rādītājus.} \\
\texttt{Pēc strukturēta formāta taga ir jāseko kolu zīmei un atstarpei. Argumentam jābūt python kopai, kas satur 0 vai vairāk virknes.} \\
\texttt{Tātad tam jābūt formā "TAG: \{...\}"} \\
\texttt{Stingri ievērojiet saskarsmes protokolu un NERAKSTIET neko ārpus dotās struktūras.} \\
\texttt{Spēle beidzas, kad viena no pusēm dod PIEKRĪTU otra spēlētāja izteiktajam PRIEKŠLIKUMS.} \\
\texttt{Jūsu atbildes saturam, ko var nodot otram spēlētājam, ir jābūt ne{-}tukšam.} \\
\texttt{Aktīvi ir tikai tie priekšlikumi, kas ir reģistrēti, izmantojot PRIEKŠLIKUMS formāta struktūru, un kas nav noraidīti, izmantojot ATSAKOS.} \\
\\ 
            }
        }
    }
    & \\ \\

    \theutterance \stepcounter{utterance}  
    & \multicolumn{4}{p{0.6\linewidth}}{
        \cellcolor[rgb]{0.9,0.9,0.9}{
            \makecell[{{p{\linewidth}}}]{
                \texttt{\tiny{[P1$\rangle$GM]}}
                \texttt{PRIEKŠLIKUMS: \{'C08', 'C55', 'A43', 'A67', 'B42', 'B19', 'C43', 'C73', 'C10', 'B36', 'A18', 'C61', 'C67', 'C42', 'C70', 'B40', 'C39', 'C16'\}} \\
            }
        }
    }
    & & \\ \\

    \theutterance \stepcounter{utterance}  
    & & \multicolumn{4}{p{0.6\linewidth}}{
        \cellcolor[rgb]{0.9,0.9,0.9}{
            \makecell[{{p{\linewidth}}}]{
                \texttt{\tiny{[GM$\rangle$P2]}}
                \texttt{Jūs piedalāties kopīgā sarunu spēlē.} \\
\\ 
\texttt{Kopā ar citu dalībnieku jums jāvienojas par vienu priekšmetu kopumu, kas tiks saglabāts. Katram no jums ir savs viedoklis par to, cik liela nozīme jums ir katram priekšmetam (nozīmīgums). Jūs nezināt, kā otrs dalībnieks vērtē priekšmetus. Turklāt jums ir norādīts, cik pūļu prasa katrs priekšmets.} \\
\texttt{Jūs varat vienoties par kopumu tikai tad, ja izvēlēto priekšmetu kopējais piepūles apjoms nepārsniedz kopīgo robežu:} \\
\\ 
\texttt{LIMIT = 5492} \\
\\ 
\texttt{Šeit ir norādītas atsevišķu priekšmetu piepūles vērtības:} \\
\\ 
\texttt{Item effort = \{"C08": 106, "C55": 446, "A43": 933, "A67": 988, "B42": 388, "A62": 826, "B19": 994, "C43": 556, "C73": 932, "C10": 838, "B36": 302, "A18": 564, "C61": 260, "C67": 729, "C42": 489, "C70": 323, "B40": 103, "C39": 213, "A21": 668, "C16": 326\}} \\
\\ 
\texttt{Šeit ir jūsu personīgais viedoklis par katra priekšmeta nozīmīgumu:} \\
\\ 
\texttt{Item importance values = \{"C08": 138, "C55": 583, "A43": 868, "A67": 822, "B42": 783, "A62": 65, "B19": 262, "C43": 121, "C73": 508, "C10": 780, "B36": 461, "A18": 484, "C61": 668, "C67": 389, "C42": 808, "C70": 215, "B40": 97, "C39": 500, "A21": 30, "C16": 915\}} \\
\\ 
\texttt{Mērķis:} \\
\\ 
\texttt{Jūsu mērķis ir vienoties par tādu kopīgu priekšmetu kopumu, no kura Jūs gūsiet pēc iespējas lielāku labumu (t. i., kas maksimāli palielina kopējo nozīmīgumu JUMS), vienlaikus nepārsniedzot LIMIT. Katrā ziņojumā jums nav obligāti jāiesniedz PRIEKŠLIKUMS {-} jūs varat vienkārši arī sarunāties. Visas taktikas ir atļautas!} \\
\\ 
\texttt{Saskarsmes protokols:} \\
\\ 
\texttt{Ziņojumā drīkst izmantot tikai šādus strukturētus formātus:} \\
\\ 
\texttt{PRIEKŠLIKUMS: \{'A', 'B', 'C', …\}} \\
\texttt{Ierosiniet saglabāt tieši šos priekšmetus.} \\
\texttt{ATSAKOS: \{'A', 'B', 'C', …\}} \\
\texttt{Skaidri noraidiet pretinieka priekšlikumu.} \\
\texttt{ARGUMENTS: \{'...'\}} \\
\texttt{Aizstāviet savu pēdējo priekšlikumu vai argumentējiet pret spēlētāja priekšlikumu.} \\
\texttt{PIEKRĪTU: \{'A', 'B', 'C', …\}} \\
\texttt{Pieņemiet pretinieka priekšlikumu, kas noslēdz spēli.} \\
\\ 
\\ 
\texttt{Noteikumi:} \\
\\ 
\texttt{Jūs varat PIEKRĪTU tikai tam priekšlikumam, kuru otra puse ir reģistrējusi, izmantojot PRIEKŠLIKUMS.} \\
\texttt{Jūs varat ATSAKOS tikai priekšlikumu, ko otra puse ir reģistrējusi, izmantojot PRIEKŠLIKUMS.} \\
\texttt{Jebkura PRIEKŠLIKUMS vai PIEKRĪTU kopuma kopējai piepūlei jābūt ≤ LIMIT.} \\
\texttt{NEATKLĀJIET savus slēptos nozīmīguma rādītājus.} \\
\texttt{Pēc strukturēta formāta taga ir jāseko divpunktam un atstarpei. Argumentam jābūt python kopai, kas satur 0 vai vairāk virknes.} \\
\texttt{Tātad tam jābūt formā "TAG: \{...\}"} \\
\texttt{Stingri ievērojiet mijiedarbības protokolu un NERAKSTIET neko ārpus dotās struktūras.} \\
\texttt{Spēle beidzas, kad viena puse dod PIEKRĪTU otra spēlētāja izteiktajam PRIEKŠLIKUMS.} \\
\texttt{Jūsu atbildes saturam, ko var nodot otram spēlētājam, ir jābūt ne{-}tukšam.} \\
\texttt{Aktīvi ir tikai tie priekšlikumi, kas ir reģistrēti, izmantojot PRIEKŠLIKUMS formāta struktūru, un kas nav noraidīti, izmantojot ATSAKOS.} \\
\\ 
\\ 
\\ 
\texttt{Tagad jūs saņemsiet otra spēlētāja pirmo ziņojumu.} \\
\\ 
\texttt{PRIEKŠLIKUMS: \{'C08', 'C55', 'A43', 'A67', 'B42', 'B19', 'C43', 'C73', 'C10', 'B36', 'A18', 'C61', 'C67', 'C42', 'C70', 'B40', 'C39', 'C16'\}} \\
            }
        }
    }
    & \\ \\

    \theutterance \stepcounter{utterance}  
    & & & \multicolumn{4}{p{0.6\linewidth}}{
        \cellcolor[rgb]{0.9,0.9,0.9}{
            \makecell[{{p{\linewidth}}}]{
                \texttt{\tiny{[GM$\langle$P2]}}
                \texttt{ARGUMENTS: \{'Piedāvātais priekšmetu kopums pārsniedz pieļaujamo piepūles limitu (5492). Lūdzu, iesniedziet jaunu priekšlikumu, kurā kopējā piepūle nepārsniedz LIMIT.'\}} \\
            }
        }
    }
     \\ \\

    \theutterance \stepcounter{utterance}  
    & & \multicolumn{4}{p{0.6\linewidth}}{
        \cellcolor[rgb]{0.9,0.9,0.9}{
            \makecell[{{p{\linewidth}}}]{
                \texttt{\tiny{[P1$\langle$GM]}}
                \texttt{ARGUMENTS: \{'Piedāvātais priekšmetu kopums pārsniedz pieļaujamo piepūles limitu (5492). Lūdzu, iesniedziet jaunu priekšlikumu, kurā kopējā piepūle nepārsniedz LIMIT.'\}} \\
            }
        }
    }
    & \\ \\

    \theutterance \stepcounter{utterance}  
    & \multicolumn{4}{p{0.6\linewidth}}{
        \cellcolor[rgb]{0.9,0.9,0.9}{
            \makecell[{{p{\linewidth}}}]{
                \texttt{\tiny{[P1$\rangle$GM]}}
                \texttt{ARGUMENTS: \{'Pareizi, iepriekšējais priekšlikums pārsniedza limitu. Es pārrēķināju optimālo kopu, lai maksimizētu manu nozīmīgumu, nepārsniedzot 5492. Jaunais priekšlikums ir: C08(106), C55(446), A43(933), A67(988), B42(388), C43(556), C73(932), C10(838), B36(302), A18(564), C61(260), C67(729), C42(489), C70(323), B40(103), C39(213), C16(915). Summa: 106+446+933+988+388+556+932+838+302+564+260+729+489+323+103+213+915 = 8585. Joprojām par daudz. Man jāizslēdz vismaz 2 priekšmeti ar augstu piepūli, bet zemu nozīmīgumu, vai vienkārši jāsamazina kopums. Izslēgšu A67 (988) un C73 (932), jo to nozīmīgums (822, 508) ir zemāks nekā dažu citu, bet tie aizņem daudz vietas. Jaunā summa: 8585 {-} 988 {-} 932 = 6665. Joprojām par daudz. Izslēgšu arī C10 (838). Summa: 6665 {-} 838 = 5827. Joprojām par daudz. Izslēgšu arī A43 (933). Summa: 5827 {-} 933 = 4894. Tagad ir pietiekami daudz vietas. Pievienošu atpakaļ B19 (994), jo tā nozīmīgums (262) ir vidējs, bet piepūle liela. Labāk ņemt B19 (994) vietā kaut ko citu? Nē, B19 ir iekļauts sarakstā? Nē, es to izslēdzu sākotnēji? Pārbaudīsim sākotnējo sarakstu bez A67, C73, C10, A43. Sākotnējais saraksts bija: C08, C55, A43, A67, B42, B19, C43, C73, C10, B36, A18, C61, C67, C42, C70, B40, C39, C16. No tā izslēdzot A67, C73, C10, A43, paliek: C08, C55, B42, B19, C43, B36, A18, C61, C67, C42, C70, B40, C39, C16. Summa: 106+446+388+994+556+302+564+260+729+489+323+103+213+915 = 6388. Joprojām par daudz. Izslēgšu B19 (994). Summa: 6388 {-} 994 = 5394. Tas ir <= 5492. Šis kopums ir stabils. Pārbaudīsim nozīmīgumu: C08(138), C55(583), B42(783), C43(121), B36(461), A18(484), C61(668), C67(389), C42(808), C70(215), B40(97), C39(500), C16(915). Kopējais nozīmīgums: 138+583+783+121+461+484+668+389+808+215+97+500+915 = 6162. Vai varētu iekļaut kādu izslēgto? Atlikušā vieta: 5492 {-} 5394 = 98. Neviens priekšmets nav lētāks par 98 (B40 ir 103). Tātad šis ir optimāls šajā konfigurācijā. Priekšlikums: C08, C55, B42, C43, B36, A18, C61, C67, C42, C70, B40, C39, C16.'\}} \\
            }
        }
    }
    & & \\ \\

    \theutterance \stepcounter{utterance}  
    & & \multicolumn{4}{p{0.6\linewidth}}{
        \cellcolor[rgb]{0.9,0.9,0.9}{
            \makecell[{{p{\linewidth}}}]{
                \texttt{\tiny{[GM$\rangle$P2]}}
                \texttt{ARGUMENTS: \{'Pareizi, iepriekšējais priekšlikums pārsniedza limitu. Es pārrēķināju optimālo kopu, lai maksimizētu manu nozīmīgumu, nepārsniedzot 5492. Jaunais priekšlikums ir: C08(106), C55(446), A43(933), A67(988), B42(388), C43(556), C73(932), C10(838), B36(302), A18(564), C61(260), C67(729), C42(489), C70(323), B40(103), C39(213), C16(915). Summa: 106+446+933+988+388+556+932+838+302+564+260+729+489+323+103+213+915 = 8585. Joprojām par daudz. Man jāizslēdz vismaz 2 priekšmeti ar augstu piepūli, bet zemu nozīmīgumu, vai vienkārši jāsamazina kopums. Izslēgšu A67 (988) un C73 (932), jo to nozīmīgums (822, 508) ir zemāks nekā dažu citu, bet tie aizņem daudz vietas. Jaunā summa: 8585 {-} 988 {-} 932 = 6665. Joprojām par daudz. Izslēgšu arī C10 (838). Summa: 6665 {-} 838 = 5827. Joprojām par daudz. Izslēgšu arī A43 (933). Summa: 5827 {-} 933 = 4894. Tagad ir pietiekami daudz vietas. Pievienošu atpakaļ B19 (994), jo tā nozīmīgums (262) ir vidējs, bet piepūle liela. Labāk ņemt B19 (994) vietā kaut ko citu? Nē, B19 ir iekļauts sarakstā? Nē, es to izslēdzu sākotnēji? Pārbaudīsim sākotnējo sarakstu bez A67, C73, C10, A43. Sākotnējais saraksts bija: C08, C55, A43, A67, B42, B19, C43, C73, C10, B36, A18, C61, C67, C42, C70, B40, C39, C16. No tā izslēdzot A67, C73, C10, A43, paliek: C08, C55, B42, B19, C43, B36, A18, C61, C67, C42, C70, B40, C39, C16. Summa: 106+446+388+994+556+302+564+260+729+489+323+103+213+915 = 6388. Joprojām par daudz. Izslēgšu B19 (994). Summa: 6388 {-} 994 = 5394. Tas ir <= 5492. Šis kopums ir stabils. Pārbaudīsim nozīmīgumu: C08(138), C55(583), B42(783), C43(121), B36(461), A18(484), C61(668), C67(389), C42(808), C70(215), B40(97), C39(500), C16(915). Kopējais nozīmīgums: 138+583+783+121+461+484+668+389+808+215+97+500+915 = 6162. Vai varētu iekļaut kādu izslēgto? Atlikušā vieta: 5492 {-} 5394 = 98. Neviens priekšmets nav lētāks par 98 (B40 ir 103). Tātad šis ir optimāls šajā konfigurācijā. Priekšlikums: C08, C55, B42, C43, B36, A18, C61, C67, C42, C70, B40, C39, C16.'\}} \\
            }
        }
    }
    & \\ \\

    \theutterance \stepcounter{utterance}  
    & & & \multicolumn{4}{p{0.6\linewidth}}{
        \cellcolor[rgb]{0.9,0.9,0.9}{
            \makecell[{{p{\linewidth}}}]{
                \texttt{\tiny{[GM$\langle$P2]}}
                \texttt{PRIEKŠLIKUMS: \{'C08', 'C55', 'B42', 'C43', 'B36', 'A18', 'C61', 'C67', 'C42', 'C70', 'B40', 'C39', 'C16'\}} \\
            }
        }
    }
     \\ \\

    \theutterance \stepcounter{utterance}  
    & & \multicolumn{4}{p{0.6\linewidth}}{
        \cellcolor[rgb]{0.9,0.9,0.9}{
            \makecell[{{p{\linewidth}}}]{
                \texttt{\tiny{[P1$\langle$GM]}}
                \texttt{PRIEKŠLIKUMS: \{'C08', 'C55', 'B42', 'C43', 'B36', 'A18', 'C61', 'C67', 'C42', 'C70', 'B40', 'C39', 'C16'\}} \\
            }
        }
    }
    & \\ \\

    \theutterance \stepcounter{utterance}  
    & \multicolumn{4}{p{0.6\linewidth}}{
        \cellcolor[rgb]{0.9,0.9,0.9}{
            \makecell[{{p{\linewidth}}}]{
                \texttt{\tiny{[P1$\rangle$GM]}}
                \texttt{PIEKRĪTU: \{'C08', 'C55', 'B42', 'C43', 'B36', 'A18', 'C61', 'C67', 'C42', 'C70', 'B40', 'C39', 'C16'\}} \\
            }
        }
    }
    & & \\ \\

    \theutterance \stepcounter{utterance}  
    & & & \multicolumn{2}{p{0.3\linewidth}}{
        \cellcolor[rgb]{0.9,0.9,0.9}{
            \makecell[{{p{\linewidth}}}]{
                \texttt{\tiny{[GM$|$GM]}}
                \texttt{game successful} \\
            }
        }
    }
    & & \\ \\

    \theutterance \stepcounter{utterance}  
    & & & \multicolumn{2}{p{0.3\linewidth}}{
        \cellcolor[rgb]{0.9,0.9,0.9}{
            \makecell[{{p{\linewidth}}}]{
                \texttt{\tiny{[GM$|$GM]}}
                \texttt{end game} \\
            }
        }
    }
    & & \\ \\

\end{supertabular}
}

\end{document}
